\section{Conclusion}
\label{sec:conclusion}

The internship began with a broad question: how much information does a model need to acquire or preserve for learned behavior to survive? The first difficulty was that ``information'' was too vague to measure directly. Much of the project therefore consisted of replacing intuitive proxies with objects that could actually be tested: LoRA rank with serialized update rate, dataset size with controlled information measures, a ``thought unit'' with boundaries scored under a stated objective, and an intermediate state with an interface tested for repeated use.

That progression also explains many of the negative results. Weight reconstruction error did not reliably predict behavioral preservation. Known source information did not determine adapter $\Rstar$. The amount by which fine-tuning improved the model was also insufficient, and the same five corpora changed their rate ordering when the frozen base model changed. On the reasoning side, no single token partition worked well for answer formation, symbolic updates, correctness, and compression. A state could be decoded without being routed into the next operation, and roughly 89\% one-step transition accuracy still collapsed after dozens of recursive uses.

Each failure made the next experiment more specific. The RMSE result moved the fine-tuning study from weight fidelity toward behavioral rate. The failure of source entropy and learning gain motivated model-relative measures of the corrections requested by the data. The disagreement between reasoning partitions shifted the project from searching for a canonical thought unit to testing explicit state interfaces. The register failure then showed that even a good interface is only useful if its update rule remains reliable under self-conditioned use.

The strongest fine-tuning result is therefore conditional rather than universal. Across the current two-model, seven-task-family discovery set, the spectral structure of the corrections requested by the training data predicts the adapter rate substantially better than source code length or model identity alone. Correction Fisher log-volume reaches mean within-model Spearman correlation 0.764, while effective rank gives the best task-held-out prediction error in the current screen. Because these measures were selected and evaluated on the same panel, they remain hypotheses for prospective replication rather than established scaling laws.

The strongest reasoning result is similarly bounded. Useful intermediate structure clearly exists, but its form depends on the future use we ask it to support. In controlled finite tasks, an explicit state interface can remain compact and reliable when it has enough capacity, one stable meaning per state, and a sufficiently accurate transition rule. The canonical four-bit proof state remains exact through long padded horizons, while lossy and path-aliased controls fail for different reasons. The mixed-register task shows the limit: modest local update errors compound until long-horizon performance is almost lost.

Taken together, these results suggest a simple discipline for future information claims. A useful rate must specify the code being transmitted, the receiver that interprets it, the side information already available to that receiver, and the behavior that must be preserved. For an intermediate state, the update rule matters as well: containing the right information once is not enough if repeated use destroys it.

The next experiments should test these ideas where the current evidence is weakest. On the fine-tuning side, content diversity should be varied while matching support and behavioral gain, and the correction-spectrum measure should be fixed before evaluating new models and tasks. On the reasoning side, one fixed reasoning adapter should be compressed while both output behavior and internal structure are tracked, and recursive state interfaces should be tested on state spaces and transition systems too large to memorize as small tables.

The report therefore does not end with one information law. It ends with a narrower and more useful picture: the amount of representation that matters is relative to a receiver, a utility, and the dynamics through which the representation will be used. Any stronger theory should reproduce the model-relative adapter result on new tasks, explain why the proof-state interface succeeds while the register interface fails, and account for the task dependence seen throughout both experimental lines.